\section{Finite modules and finitely presented modules}
\label{section-module-finite-type}

\noindent
Just some basic notation and lemmas.

\begin{definition}
\label{definition-module-finite-type}
Let $R$ be a ring. Let $M$ be an $R$-module.
\begin{enumerate}
\item We say $M$ is a {\it finite $R$-module}, or a {\it finitely generated
$R$-module} if there exist $n \in \mathbf{N}$ and $x_1, \ldots, x_n \in M$
such that every element of $M$ is an $R$-linear combination of the $x_i$.
Equivalently, this means there exists a surjection
$R^{\oplus n} \to M$ for some $n \in \mathbf{N}$.
\item We say $M$ is a {\it finitely presented $R$-module} or an
{\it $R$-module of finite presentation} if there exist integers
$n, m \in \mathbf{N}$ and an exact sequence
$$
R^{\oplus m} \longrightarrow R^{\oplus n} \longrightarrow M \longrightarrow 0
$$
\end{enumerate}
\end{definition}

\noindent
Informally, $M$ is a finitely presented $R$-module if and only if
it is finitely generated and the module of relations among these
generators is finitely generated as well.
A choice of an exact sequence as in the definition is called a
{\it presentation} of $M$.

\begin{lemma}
\label{lemma-lift-map}
Let $R$ be a ring. Let $\alpha : R^{\oplus n} \to M$ and $\beta : N \to M$ be
module maps. If $\Im(\alpha) \subset \Im(\beta)$, then there
exists an $R$-module map $\gamma : R^{\oplus n} \to N$ such that
$\alpha = \beta \circ \gamma$.
\end{lemma}

\begin{proof}
Let $e_i = (0, \ldots, 0, 1, 0, \ldots, 0)$ be the $i$th basis vector
of $R^{\oplus n}$. Let $x_i \in N$ be an element with
$\alpha(e_i) = \beta(x_i)$ which exists by assumption. Set
$\gamma(a_1, \ldots, a_n) = \sum a_i x_i$. By construction
$\alpha = \beta \circ \gamma$.
\end{proof}

\begin{lemma}
\label{lemma-extension}
Let $R$ be a ring.
Let
$$
0 \to M_1 \to M_2 \to M_3 \to 0
$$
be a short exact sequence of $R$-modules.
\begin{enumerate}
\item If $M_1$ and $M_3$ are finite $R$-modules, then $M_2$ is a finite
$R$-module.
\item If $M_1$ and $M_3$ are finitely presented $R$-modules, then $M_2$
is a finitely presented $R$-module.
\item If $M_2$ is a finite $R$-module, then $M_3$ is a finite $R$-module.
\item If $M_2$ is a finitely presented $R$-module and $M_1$ is a
finite $R$-module, then $M_3$ is a finitely presented $R$-module.
\item If $M_3$ is a finitely presented $R$-module and $M_2$ is a finite
$R$-module, then $M_1$ is a finite $R$-module.
\end{enumerate}
\end{lemma}

\begin{proof}
Proof of (1). If $x_1, \ldots, x_n$ are generators of $M_1$ and
$y_1, \ldots, y_m \in M_2$ are elements whose images in $M_3$ are
generators of $M_3$, then $x_1, \ldots, x_n, y_1, \ldots, y_m$
generate $M_2$.

\medskip\noindent
Part (3) is immediate from the definition.

\medskip\noindent
Proof of (5). Assume $M_3$ is finitely presented and $M_2$ finite.
Choose a presentation
$$
R^{\oplus m} \to R^{\oplus n} \to M_3 \to 0
$$
By Lemma \ref{lemma-lift-map} there exists a map
$R^{\oplus n} \to M_2$ such that
the solid diagram
$$
\xymatrix{
& R^{\oplus m} \ar[r] \ar@{..>}[d] & R^{\oplus n} \ar[r] \ar[d] &
M_3 \ar[r] \ar[d]^{\text{id}} & 0 \\
0 \ar[r] & M_1 \ar[r] & M_2 \ar[r] & M_3 \ar[r] & 0
}
$$
commutes. This produces the dotted arrow. By the snake lemma
(Lemma \ref{lemma-snake}) we see that we get an isomorphism
$$
\Coker(R^{\oplus m} \to M_1)
\cong
\Coker(R^{\oplus n} \to M_2)
$$
In particular we conclude that $\Coker(R^{\oplus m} \to M_1)$
is a finite $R$-module. Since $\Im(R^{\oplus m} \to M_1)$
is finite by (3), we see that $M_1$ is finite by part (1).

\medskip\noindent
Proof of (4). Assume $M_2$ is finitely presented and $M_1$ is finite.
Choose a presentation $R^{\oplus m} \to R^{\oplus n} \to M_2 \to 0$.
Choose a surjection $R^{\oplus k} \to M_1$. By Lemma \ref{lemma-lift-map}
there exists a factorization $R^{\oplus k} \to R^{\oplus n} \to M_2$
of the composition $R^{\oplus k} \to M_1 \to M_2$. Then
$R^{\oplus k + m} \to R^{\oplus n} \to M_3 \to 0$
is a presentation.

\medskip\noindent
Proof of (2). Assume that $M_1$ and $M_3$ are finitely presented.
The argument in the proof of part (1) produces a commutative diagram
$$
\xymatrix{
0 \ar[r] & R^{\oplus n} \ar[d] \ar[r] & R^{\oplus n + m} \ar[d] \ar[r] &
R^{\oplus m} \ar[d] \ar[r] & 0 \\
0 \ar[r] & M_1 \ar[r] & M_2 \ar[r] & M_3 \ar[r] & 0
}
$$
with surjective vertical arrows. By the snake lemma we obtain a short
exact sequence
$$
0 \to \Ker(R^{\oplus n} \to M_1) \to
\Ker(R^{\oplus n + m} \to M_2) \to
\Ker(R^{\oplus m} \to M_3) \to 0
$$
By part (5) we see that the outer two modules are finite. Hence the
middle one is finite too. By (4) we see that $M_2$ is of finite presentation.
\end{proof}

\begin{lemma}
\label{lemma-trivial-filter-finite-module}
\begin{slogan}
Finite modules have filtrations such that successive quotients are
cyclic modules.
\end{slogan}
Let $R$ be a ring, and let $M$ be a finite $R$-module.
There exists a filtration by finite $R$-submodules
$$
0 = M_0 \subset M_1 \subset \ldots \subset M_n = M
$$
such that each quotient $M_i/M_{i - 1}$ is isomorphic
to $R/I_i$ for some ideal $I_i$ of $R$.
\end{lemma}

\begin{proof}
By induction on the number of generators of $M$. Let
$x_1, \ldots, x_r \in M$ be generators.
Let $M' = Rx_1 \subset M$. Then $M/M'$ has $r - 1$ generators
and the induction hypothesis applies. And clearly $M' \cong R/I_1$
with $I_1 = \{f \in R \mid fx_1 = 0\}$.
\end{proof}

\begin{lemma}
\label{lemma-finite-over-subring}
Let $R \to S$ be a ring map.
Let $M$ be an $S$-module.
If $M$ is finite as an $R$-module, then $M$ is finite as an $S$-module.
\end{lemma}

\begin{proof}
In fact, any $R$-generating set of $M$ is also an $S$-generating set of
$M$, since the $R$-module structure is induced by the image of $R$ in $S$.
\end{proof}
